\begin{table}[H]\centering\footnotesize
\caption{End-to-end learning of the geometry fails (Scenario 4).}
\label{tab:collapse}
\setlength{\tabcolsep}{5pt}
\resizebox{\columnwidth}{!}{%
\begin{tabular}{@{}l r r r@{}}
\toprule
Dataset & fixed+random & fixed+learned wts & \textbf{learned geom.}\\
\midrule
MSRC-v5      & 0.682 & 0.470 & 0.328\\
Caltech101-7 & 0.535 & 0.410 & \textbf{0.049}\\
OutdoorScene & 0.570 & 0.463 & \textbf{0.005}\\
UCI          & 0.822 & 0.914 & 0.751\\
\bottomrule
\end{tabular}}
\end{table}

\noindent\textbf{Setup.} AMI as progressively more of the operator is made trainable (left$\to$right).
fixed+random: the per-view kernels $P_v$ (Gaussian kNN affinities) and the fusion are both fixed,
a random convex trajectory, no learning. fixed+learned wts: keep the kernels $P_v$ fixed and
learn only the mixing weights, the $t\times V$ matrix $a$ (one softmax over views per step) that
convexly combines the fixed per-view walks, $\W=\prod_s\big(\sum_v a_{sv}P_v\big)$, by Adam on the
contrastive loss (positive pairs $=$ kernel neighbours). learned geom.: also learn the geometry
itself: a per-view feature map $\phi_v$ defines each kernel $P_v(\phi_v)$, so both the kernels
and $a$ are trained end-to-end under the same loss (now even the Gram target moves). Four
datasets; bold $=$ near-zero collapse.

\noindent\textbf{Result.} Accuracy falls monotonically as more is learned: fixed $>$ learned-weights $>$
fully-learned (near-zero on Caltech/OutdoorScene). Learning the operator hurts. The mechanism is
optimisation/conditioning (flat loss, spectrum flattening), not the literal rank-one collapse the
numbers suggest.
